\input{../tex-inputs/latex-leseansicht-vorspann}

\section[Gustav und Max Schwarzkopf an Arthur Schnitzler, 11. 5. 1904]{L04313 Gustav und Max Schwarzkopf an Arthur Schnitzler, 11. 5. 1904}
\nopagebreak\mylabel{L04313v}
\rehead{ }\normalsize\beginnumbering\briefempfaengerindex{Schnitzler, Arthur@\textsc{Schnitzler, Arthur}!zzzSchwarzkopf, Max@\emph{von Max Schwarzkopf}!1904-05-111@{11. 5. 1904}|(be}\briefempfaengerindex{Schnitzler, Arthur@\textsc{Schnitzler, Arthur}!zzzSchwarzkopf, Gustav@\emph{von Gustav Schwarzkopf}!1904-05-111@{11. 5. 1904}|(be}
\toendnotes[C]{\smallbreak\pagebreak[2]}
\correspDesc{Versand  durch Gustav Schwarzkopf, Max Schwarzkopf am 11. 5. 1904 in Wien
\newline{}Erhalt  durch Arthur Schnitzler im Zeitraum [12. 5. 1904
                  – 16. 5. 1904?] in Wien}\toendnotes[C]{\smallbreak}
\Standort{CUL, Schnitzler, B 96a.}
\physDesc{Brief, 1 Blatt, 4 Seiten, 2277 Zeichen
\newline{}Handschrift: schwarze Tinte, deutsche Kurrent}\toendnotes[C]{\smallbreak}
\pstart
           \raggedleft{}{\pb}11. V. 04\pend
           \vspace{0.5em}
\pstart
           Lieber Arthur, ich habe ihn gefunden, den »reinen Tor\pwindex{Schwarzkopf, Max 12.\,6.\,1857 Wien – 14.\,4.\,1928 ebd.@\textsc{Schwarzkopf, Max} (12.\,6.\,1857 Wien – 14.\,4.\,1928 ebd.), \emph{Rechtsanwalt}!reine Tor. Gesellschaftsstück in vier Akten@\strich\emph{Der reine Tor. Gesellschaftsstück in vier Akten}|pw}«, er war in dem kleinen Schrank. Bei der
               Gelenheit habe ich den wirklichen reinen Tor, Ihren Heini\pwindex{Schnitzler, Heinrich 9.\,8.\,1902 Hinterbrühl – 12.\,7.\,1982 Wien@\textsc{Schnitzler, Heinrich} (9.\,8.\,1902 Hinterbrühl – 12.\,7.\,1982 Wien), \emph{Regisseur, Schauspieler}|pw} geſehen. (Ihre Frau\pwindex{Schnitzler, Olga 17.\,1.\,1882 Wien – 13.\,1.\,1970 Lugano@\textsc{Schnitzler, Olga} (17.\,1.\,1882 Wien – 13.\,1.\,1970 Lugano), \emph{Schauspielerin, Sängerin}|pwv} wird dieſe Bene{\geminationn}ung hoffentlich nicht als
               Herabſetzung der geiſtigen Fähigkeiten Heini’s\pwindex{Schnitzler, Heinrich 9.\,8.\,1902 Hinterbrühl – 12.\,7.\,1982 Wien@\textsc{Schnitzler, Heinrich} (9.\,8.\,1902 Hinterbrühl – 12.\,7.\,1982 Wien), \emph{Regisseur, Schauspieler}|pw}
               empfinden) Er hat mich sehr gnädig empfangen, hat mir alle{ }ſeine Kunſtſtücke
               vorgemacht, hat unter Muſikbegleitung bereitwilligſt vor mir getanzt, ohne auch nur
               den allergeringſten Kopf dafür zu verlangen. Da{\geminationn} hat er
               mir {\pb}erzählt, daß er im »Ratauspak\oindex{Wien@\textbf{Wien}!I., Innere Stadt@\textbf{I., Innere Stadt}!Rathauspark@\textbf{Rathauspark}, \emph{Park}|pw}« den »Bi{\geminationn}gbu{\geminationn}en« gesehen hat und daß er eine »Snitzler
               Locke« hat. (Man hat ihm wirklich eine gemacht,{ }ſeine Mittel erlauben ihm das
               bereits) Auch über{ }ſeinen Geſundheitszuſtand kö{\geminationn}en Sie
               beruhigt{ }ſein, de{\geminationn} er hat, gleichſam als Glanznu{\geminationm}er und letzte »Production« in meiner Gegenwart gezeigt,
               daß alles tadellos funktioniert und daß er nicht an Verdauungsſtörungen leidet.\pend
           
\pstart
           Kurz, ich habe mich mit ihm beſſer unterhalten als \label{K_L04313-1v}\edtext{neulich im »Carl Theater\oindex{Wien@\textbf{Wien}!II., Leopoldstadt@\textbf{II., Leopoldstadt}!Carl-Theater@\textbf{Carl-Theater}, \emph{Theater}|pw}\eventindex{Carl-Theater@\textbf{Carl-Theater}!Premiere von Das Tal des Lebens, 6.5.1904@Premiere von Das Tal des Lebens, 6.5.1904|pwv}« bei »Tal des Lebens\pwindex{Dreyer, Max 25.\,9.\,1862 Rostock – 27.\,11.\,1946 Göhren@\textsc{Dreyer, Max} (25.\,9.\,1862 Rostock – 27.\,11.\,1946 Göhren), \emph{Schriftsteller}!Tal des Lebens. Historischer Schwank in 4 Aufzügen@\strich\emph{Das Tal des Lebens. Historischer Schwank in 4 Aufzügen}|pw}}{\lemma{\textnormal{\emph{neulich … Lebens}}}\Cendnote{\textnormal{Vom 3. 5. 1904
                  bis zum 31. 5. 1904 fand ein »Gesamtgastspiel« des \emph{Deutschen Theaters Berlin}\orgindex{Deutsches Theater Berlin@Deutsches Theater Berlin|pwk} im Carl-Theater\oindex{Wien@\textbf{Wien}!II., Leopoldstadt@\textbf{II., Leopoldstadt}!Carl-Theater@\textbf{Carl-Theater}, \emph{Theater}|pwk} statt. Gegeben wurden (mehrfach) drei Stücke:
                     \emph{Die Weber}\pwindex{Hauptmann, Gerhart 15.\,11.\,1862 Szczawno-Zdrój – 6.\,6.\,1946 Jagniątków@\textsc{Hauptmann, Gerhart} (15.\,11.\,1862 Szczawno-Zdrój – 6.\,6.\,1946 Jagniątków), \emph{Schriftsteller}!Weber. Schauspiel aus den vierziger Jahren@\strich\emph{Die Weber. Schauspiel aus den vierziger Jahren}|pwk} (von Gerhart Hauptmann\pwindex{Hauptmann, Gerhart 15.\,11.\,1862 Szczawno-Zdrój – 6.\,6.\,1946 Jagniątków@\textsc{Hauptmann, Gerhart} (15.\,11.\,1862 Szczawno-Zdrój – 6.\,6.\,1946 Jagniątków), \emph{Schriftsteller}|pwk}), \emph{Das
                     Tal des Lebens}\pwindex{Dreyer, Max 25.\,9.\,1862 Rostock – 27.\,11.\,1946 Göhren@\textsc{Dreyer, Max} (25.\,9.\,1862 Rostock – 27.\,11.\,1946 Göhren), \emph{Schriftsteller}!Tal des Lebens. Historischer Schwank in 4 Aufzügen@\strich\emph{Das Tal des Lebens. Historischer Schwank in 4 Aufzügen}|pwk} (von Max Dreyer\pwindex{Dreyer, Max 25.\,9.\,1862 Rostock – 27.\,11.\,1946 Göhren@\textsc{Dreyer, Max} (25.\,9.\,1862 Rostock – 27.\,11.\,1946 Göhren), \emph{Schriftsteller}|pwk}) 
                  und \emph{Der Meister}\pwindex{Bahr, Hermann 19.\,7.\,1863 Linz – 15.\,1.\,1934 München@\textsc{Bahr, Hermann} (19.\,7.\,1863 Linz – 15.\,1.\,1934 München), \emph{Schriftsteller, Kritiker}!Meister. Komödie in drei Akten@\strich\emph{Der Meister. Komödie in drei Akten}|pwk} (von Hermann Bahr\pwindex{Bahr, Hermann 19.\,7.\,1863 Linz – 15.\,1.\,1934 München@\textsc{Bahr, Hermann} (19.\,7.\,1863 Linz – 15.\,1.\,1934 München), \emph{Schriftsteller, Kritiker}|pwk}). Die Premiere von \emph{Das Tal des Lebens}\pwindex{Dreyer, Max 25.\,9.\,1862 Rostock – 27.\,11.\,1946 Göhren@\textsc{Dreyer, Max} (25.\,9.\,1862 Rostock – 27.\,11.\,1946 Göhren), \emph{Schriftsteller}!Tal des Lebens. Historischer Schwank in 4 Aufzügen@\strich\emph{Das Tal des Lebens. Historischer Schwank in 4 Aufzügen}|pwk}\eventindex{Carl-Theater@\textbf{Carl-Theater}!Premiere von Das Tal des Lebens, 6.5.1904@Premiere von Das Tal des Lebens, 6.5.1904|pwk}
                  fand am 6. 5. 1904 statt. }}}\label{K_L04313-1}«. Es war nicht
               meine {\pb}Wahl, ich wurde mitgenommen.
               Das Stück\pwindex{Dreyer, Max 25.\,9.\,1862 Rostock – 27.\,11.\,1946 Göhren@\textsc{Dreyer, Max} (25.\,9.\,1862 Rostock – 27.\,11.\,1946 Göhren), \emph{Schriftsteller}!Tal des Lebens. Historischer Schwank in 4 Aufzügen@\strich\emph{Das Tal des Lebens. Historischer Schwank in 4 Aufzügen}|pwv} iſt eine nette
               Anekdote, die für einen Act reichen würde, und die mit primitiver Technik zu vier
               Acten zerdehnt wird. Hie und da ein hübsches Witzwort, aber alles wird zehnmal geſagt
               und jede Zweideutigkeit{ }ſo gehetzt und breitgetreten, bis{ }ſie den letzten Reſt von
               Pikanterie verliert. Darſtellung (bis auf die Trieſch\pwindex{Triesch, Irene 13.\,4.\,1877 Wien – 24.\,11.\,1964 Basel@\textsc{Triesch, Irene} (13.\,4.\,1877 Wien – 24.\,11.\,1964 Basel), \emph{Schauspielerin}|pw}) Ausſtattung, Regie unwahrſcheinlich{ }ſchlecht. Es iſt auch ganz
               abgefallen, nur der dritte Act\pwindex{Dreyer, Max 25.\,9.\,1862 Rostock – 27.\,11.\,1946 Göhren@\textsc{Dreyer, Max} (25.\,9.\,1862 Rostock – 27.\,11.\,1946 Göhren), \emph{Schriftsteller}!Tal des Lebens. Historischer Schwank in 4 Aufzügen@\strich\emph{Das Tal des Lebens. Historischer Schwank in 4 Aufzügen}|pwv} hatte Beifall. Die Kritik iſt diesmal dem Gaſtſpiel\orgindex{Deutsches Theater Berlin@Deutsches Theater Berlin|pwv} gegenüber recht kritiſch, das Publikum{ }ſehr kühl,
               aber gute Häuſer machen{ }ſie doch, de{\geminationn} es iſt kalt und
               unfreundlich, der Frühling iſt wieder einmal {\pb}ein »grünangeſtrichener Winter«, ein
               richtiges Brahm\pwindex{Brahm, Otto 5.\,2.\,1856 Hamburg – 28.\,11.\,1912 Berlin@\textsc{Brahm, Otto} (5.\,2.\,1856 Hamburg – 28.\,11.\,1912 Berlin), \emph{Theaterleiter, Regisseur}|pw}-Wetter\substVorne{}\textsuperscript{!}\substDazwischen{}.\substHinten{} Ob er aber mit{ }ſeinen drei Stücken\pwindex{Hauptmann, Gerhart 15.\,11.\,1862 Szczawno-Zdrój – 6.\,6.\,1946 Jagniątków@\textsc{Hauptmann, Gerhart} (15.\,11.\,1862 Szczawno-Zdrój – 6.\,6.\,1946 Jagniątków), \emph{Schriftsteller}!Weber. Schauspiel aus den vierziger Jahren@\strich\emph{Die Weber. Schauspiel aus den vierziger Jahren}|pwv}\pwindex{Bahr, Hermann 19.\,7.\,1863 Linz – 15.\,1.\,1934 München@\textsc{Bahr, Hermann} (19.\,7.\,1863 Linz – 15.\,1.\,1934 München), \emph{Schriftsteller, Kritiker}!Meister. Komödie in drei Akten@\strich\emph{Der Meister. Komödie in drei Akten}|pwv}\pwindex{Dreyer, Max 25.\,9.\,1862 Rostock – 27.\,11.\,1946 Göhren@\textsc{Dreyer, Max} (25.\,9.\,1862 Rostock – 27.\,11.\,1946 Göhren), \emph{Schriftsteller}!Tal des Lebens. Historischer Schwank in 4 Aufzügen@\strich\emph{Das Tal des Lebens. Historischer Schwank in 4 Aufzügen}|pwv} doch das
               Auslangen finden wird? Vielleicht ko{\geminationm}en Sie doch noch
               daran. – \textsc{Angelo\pwindex{Neumann, Angelo 18.\,8.\,1838 Stupava – 20.\,12.\,1910 Prag@\textsc{Neumann, Angelo} (18.\,8.\,1838 Stupava – 20.\,12.\,1910 Prag), \emph{Theaterleiter, Sänger}|pw}} tut{ }ſo, als ob er wirklich Ernſt machen wollte, er hat{ }ſo dringend die beiden
               Exemplare verlangt als ob die Proben\pwindex{Schwarzkopf, Max 12.\,6.\,1857 Wien – 14.\,4.\,1928 ebd.@\textsc{Schwarzkopf, Max} (12.\,6.\,1857 Wien – 14.\,4.\,1928 ebd.), \emph{Rechtsanwalt}!reine Tor. Gesellschaftsstück in vier Akten@\strich\emph{Der reine Tor. Gesellschaftsstück in vier Akten}|pwv} demnächſt begi{\geminationn}en{ }ſollten.–  \label{K_L04313-2v}\edtext{Hofmannsthal\pwindex{Hofmannsthal, Hugo von 1.\,2.\,1874 Wien – 15.\,7.\,1929 Rodaun@\textsc{Hofmannsthal, Hugo von} (1.\,2.\,1874 Wien – 15.\,7.\,1929 Rodaun), \emph{Schriftsteller}|pw}\pwindex{Hofmannsthal, Gertrude von 16.\,3.\,1880 Wien – 9.\,11.\,1959 Paddington@\textsc{Hofmannsthal, Gertrude von} (16.\,3.\,1880 Wien – 9.\,11.\,1959 Paddington)|pw} reiſen am 18.}{\lemma{\textnormal{\emph{Hofmannsthal … 18.}}}\Cendnote{\textnormal{Hugo\pwindex{Hofmannsthal, Hugo von 1.\,2.\,1874 Wien – 15.\,7.\,1929 Rodaun@\textsc{Hofmannsthal, Hugo von} (1.\,2.\,1874 Wien – 15.\,7.\,1929 Rodaun), \emph{Schriftsteller}|pwk}, Gerty\pwindex{Hofmannsthal, Gertrude von 16.\,3.\,1880 Wien – 9.\,11.\,1959 Paddington@\textsc{Hofmannsthal, Gertrude von} (16.\,3.\,1880 Wien – 9.\,11.\,1959 Paddington)|pwk} und sein Vater\pwindex{Hofmannsthal, Hugo August von 21.\,12.\,1841 Wien – 8.\,12.\,1915 ebd.@\textsc{Hofmannsthal, Hugo August von} (21.\,12.\,1841 Wien – 8.\,12.\,1915 ebd.), \emph{Bankdirektor}|pwkv} reisten nach Holland\oindex{Niederlande@\textbf{Niederlande}|pwk}. }}}\label{K_L04313-2}
               fort, wie mir der Page geſtern erzählte. –\pend
           
\pstart
           Daß ich nichts Mitteilenswertes erlebt habe, glauben Sie mir wol, ich bezweifle aber
               doch nicht, daß die Welt an{ }ſich{ }ſchön iſt.\pend
           
\pstart
           Ihre Frau\pwindex{Schnitzler, Olga 17.\,1.\,1882 Wien – 13.\,1.\,1970 Lugano@\textsc{Schnitzler, Olga} (17.\,1.\,1882 Wien – 13.\,1.\,1970 Lugano), \emph{Schauspielerin, Sängerin}|pwv} und Sie
               herzlichſt grüßend,{\\[\baselineskip]}Ihr \spacefill\mbox{Gustav}\pend
           \leftskip=0em{}\selectlanguage{ngerman}\vspace{1em}{\vspace{1\baselineskip}}
\pstart
           {[}hs. Schwarzkopf:{]} Natürlich habe ich Grüße von \textsc{Heini\pwindex{Schnitzler, Heinrich 9.\,8.\,1902 Hinterbrühl – 12.\,7.\,1982 Wien@\textsc{Schnitzler, Heinrich} (9.\,8.\,1902 Hinterbrühl – 12.\,7.\,1982 Wien), \emph{Regisseur, Schauspieler}|pw}} zu beſtellen\pend
           \pstart Sie und Ihre Frau\pwindex{Schnitzler, Olga 17.\,1.\,1882 Wien – 13.\,1.\,1970 Lugano@\textsc{Schnitzler, Olga} (17.\,1.\,1882 Wien – 13.\,1.\,1970 Lugano), \emph{Schauspielerin, Sängerin}|pwv} grüßt
               herzlichſt \spacefill\mbox{Max.}\pend{}\selectlanguage{ngerman}\endnumbering\briefempfaengerindex{Schnitzler, Arthur@\textsc{Schnitzler, Arthur}!zzzSchwarzkopf, Max@\emph{von Max Schwarzkopf}!1904-05-111@{11. 5. 1904}|)be}\briefempfaengerindex{Schnitzler, Arthur@\textsc{Schnitzler, Arthur}!zzzSchwarzkopf, Gustav@\emph{von Gustav Schwarzkopf}!1904-05-111@{11. 5. 1904}|)be}\mylabel{L04313h}
\begin{anhang}
\end{anhang}\newcommand{\dateiname}{L04313}\newcommand{\titel}{Gustav und Max Schwarzkopf an Arthur Schnitzler, 11. 5. 1904}\newcommand{\editorInnen}{Herausgegeben von Jahnke, SelmaMüller, Martin Anton}\input{../tex-inputs/latex-leseansicht-abspann}
